\documentclass[sigconf]{acmart}

\settopmatter{printacmref=false}
\setcopyright{none}
\renewcommand\footnotetextcopyrightpermission[1]{}

\usepackage{booktabs}
\usepackage{amsmath}
\usepackage{listings}
\usepackage[T1]{fontenc}
\usepackage[utf8]{inputenc}
\usepackage{xcolor}
\usepackage{microtype}

\lstset{
  basicstyle=\ttfamily\footnotesize,
  breaklines=true,
  frame=single,
  backgroundcolor=\color{gray!8}
}

\title{A Disk-Backed Search Engine for Project Gutenberg:\\
Structured Retrieval, TF-IDF Cosine Similarity, and BM25}

\author{Rameez Shafat}
\affiliation{%
  \institution{Dublin City University}
  \city{Dublin}
  \country{Ireland}
}
\email{rameez.shafat2@mail.dcu.ie}

\begin{document}

\begin{abstract}
This report presents the design and implementation of an information retrieval (IR) system built over a large subset of the Project Gutenberg public-domain book collection. The system supports three retrieval approaches: structured metadata search, a Vector Space Model (VSM) using TF-IDF cosine similarity, and BM25 probabilistic ranking. To handle corpus scale safely without memory exhaustion, indexing is implemented as a disk-backed batch pipeline that writes partial inverted indexes to SQLite databases, merges them into a unified index, and precomputes collection statistics used at query time. Evaluation is performed over six fixed assignment queries, producing top-100 ranked result lists per model exported as TSV files, alongside comparative metrics including Intersection@10, Intersection@100, and ranking displacement. Results demonstrate that TF-IDF and BM25 retrieve strongly overlapping candidate sets but differ meaningfully in ranking order, while structured retrieval exhibits an orthogonal retrieval profile because it operates exclusively over bibliographic metadata fields. The report discusses design trade-offs, quantitative evaluation findings, and directions for future improvement.
\end{abstract}

\maketitle

% ============================================================
\section{Introduction}
% ============================================================

Information retrieval systems must locate relevant documents efficiently from large, heterogeneous collections. A naive linear scan of every document for every query is computationally prohibitive at scale: even with modern hardware, scanning thousands of full-text books per query produces unacceptable latency. The well-established solution is to build an inverted index offline so that query-time work is proportional to the size of postings lists for query terms rather than to the total corpus size \cite{manning2008,zobel2006}.

The goal of this project is to build a complete retrieval pipeline over a copy of the Project Gutenberg corpus. The corpus consists of plain-text book files accompanied by a metadata CSV containing bibliographic fields such as title, author, language, and bookshelf category. The system must satisfy three hard constraints. First, it must remain memory-stable during indexing despite the large corpus volume. Second, it must implement all retrieval logic directly in code without delegating to external search engines or IR frameworks. Third, it must support a controlled, reproducible comparison of three distinct retrieval approaches on a fixed set of six queries.

The architecture cleanly separates offline work from online work. During the offline phase, documents are processed in batches, written as partial index files, merged into a final index, and used to precompute statistics required by each ranking model. During the online phase, only query-relevant postings are fetched and scored. This build-time/query-time separation is a standard design principle that minimises query latency by front-loading expensive computation \cite{manning2008}.

\subsection{System Architecture}

The implementation is decomposed into six modules, each with a single well-defined responsibility:

\begin{enumerate}
    \item \textbf{Dataset loader:} parses the metadata CSV, validates that each entry has a corresponding plain-text file, and filters duplicates on \texttt{gutenberg\_id}.
    \item \textbf{Preprocessor:} normalises text by lowercasing and tokenises using a Unicode-aware alphanumeric regular expression, with optional stopword removal and Porter stemming.
    \item \textbf{Batch indexer:} processes documents in configurable batches, writing one partial SQLite index per batch before merging all partials into a final index.
    \item \textbf{Statistics module:} computes and persists document frequencies (\texttt{df}), document lengths, average document length (\texttt{avgdl}), and TF-IDF document norms.
    \item \textbf{Retriever:} implements the three scoring models and returns ranked result lists.
    \item \textbf{Evaluator:} runs the six hardcoded assignment queries over all three models, exports TSV result files, and writes a comparative summary.
\end{enumerate}

A command-line interface exposes three top-level commands: \texttt{build-index} for offline preparation, \texttt{search} for interactive single-query retrieval, and \texttt{run-experiments} for the full reproducible evaluation run.

% ============================================================
\section{Indexing}
% ============================================================

\subsection{Document Pre-processing}

Each plain-text document undergoes a normalisation pipeline before indexing: (1) decode the file using UTF-8 with error replacement to handle encoding anomalies common in older Project Gutenberg texts; (2) lowercase all characters to collapse case variants; (3) tokenise using the Unicode-aware pattern \texttt{\textbackslash w+}, which captures alphanumeric sequences and preserves accented characters; (4) discard tokens shorter than two characters to remove noise symbols.

This baseline approach was chosen deliberately over more aggressive preprocessing for two reasons. First, Manning et al.\ demonstrate that simple tokenisation is often sufficient for English-language corpora when the goal is lexical matching, and that over-normalisation can harm precision by conflating semantically distinct terms \cite{manning2008}. Second, the Project Gutenberg corpus contains texts in multiple languages --- for example, the query \emph{Dornr\"{o}schen} targets German-language books --- and applying English-specific stemming or stopword lists uniformly would damage retrieval quality for non-English documents without providing a commensurate benefit.

Two optional preprocessing switches are supported and must be activated at index build time:
\begin{itemize}
    \item \texttt{--use-stopwords}: removes common English stopwords using the NLTK list if available, falling back to a hardcoded list of 150 high-frequency English function words otherwise.
    \item \texttt{--use-stemming}: applies the Porter stemmer \cite{manning2008} to all tokens, reducing morphological variants to a shared stem.
\end{itemize}

The active configuration is written to \texttt{processed/stats/preprocessing\_config.json} at build time, and query-time tokenisation reads this same file. This guarantees that index and query apply identical preprocessing --- a common source of subtle retrieval bugs if left unmanaged.

Metadata is validated before indexing begins. The loader rejects rows with missing \texttt{gutenberg\_id} values, removes duplicate entries, and confirms that a corresponding \texttt{books/<id>.txt} file exists on disk, preventing silent index corruption from broken metadata references.

\subsection{Inverted Index Construction}

The core data structure is an inverted index backed by SQLite with the following schema:

\begin{itemize}
    \item \texttt{postings(term, doc\_id, tf)} --- stores term-document frequency triples, indexed on \texttt{term} for efficient postings lookup.
    \item \texttt{documents(doc\_id, length)} --- stores per-document token counts required for BM25 length normalisation.
    \item \texttt{metadata(doc\_id, title, author, language, bookshelf)} --- stores bibliographic fields used by structured retrieval.
\end{itemize}

SQLite was selected over alternatives such as plain JSON postings files or in-memory dictionaries for three reasons: it provides disk-backed persistence with transactional safety, supports efficient equality and range lookups via B-tree indexes, and requires no separate server process. Because the assignment prohibits external IR frameworks, SQLite represents a practical middle ground between engineering effort and runtime performance \cite{zobel2006}.

\subsection{Batch Processing and Memory Safety}

Corpus scale makes single-pass in-memory indexing unsafe. The batch indexer processes documents in configurable-sized groups (default: 250 per batch). For each batch: (1) a subset of document files is read and tokenised; (2) in-memory term-frequency maps are built for that subset only; (3) a partial index file is written to disk; (4) all temporary in-memory structures are discarded before the next batch begins.

After all batches complete, the partial indexes are merged into the final database using SQL \texttt{ATTACH DATABASE} and \texttt{INSERT INTO \ldots SELECT} statements, delegating bulk data movement to SQLite's internal engine rather than Python row iteration. The postings index is dropped before bulk insertion and recreated afterwards, avoiding the per-row index-maintenance overhead that would otherwise make the merge prohibitively slow on large corpora.

\subsection{Operational Robustness}

Two safeguards address the vulnerability of long-running builds to interruption:

\begin{enumerate}
    \item \textbf{Resumable checkpointing:} after each successful batch, its number is appended to \texttt{index\_checkpoint.json}. Rerunning with \texttt{--resume} skips already-completed batches, allowing builds to continue from the failure point.
    \item \textbf{Disk space monitoring:} available disk space is checked before each new partial index is created. If free space falls below a configurable threshold, the build raises a readable error rather than failing with a cryptic OS exception.
\end{enumerate}

A \texttt{--dry-run} mode samples a representative document subset, estimates total token volume, projected index size, and approximate build time, then exits without writing any files --- useful for capacity planning on constrained infrastructure.

% ============================================================
\section{Ranking}
% ============================================================

\subsection{Structured Metadata Retrieval}

Structured retrieval matches query terms against the bibliographic metadata fields stored in the \texttt{metadata} table. Each term match contributes a field-specific weight to the document score:

\begin{itemize}
    \item \textbf{title} = 3.0 --- the title is the primary topical descriptor of a book and carries the most concentrated relevance signal for user intent.
    \item \textbf{author} = 2.0 --- author name is highly informative for attribution queries (e.g., \emph{Philip K Dick}) but is less discriminative than title for thematic queries.
    \item \textbf{bookshelf} = 1.0 --- bookshelf is a coarse category label that provides supplementary signal but is too broad to serve as a primary ranking criterion.
\end{itemize}

These weights reflect a deliberate priority ordering based on information density and correlation with user intent, following standard practice in fielded retrieval \cite{manning2008}. All weights are configurable via the CLI for controlled ablation experiments.

Structured retrieval is fast and highly interpretable for metadata-centric queries. Its fundamental limitation is that it produces no results for queries where target terms do not appear in any metadata field, regardless of how frequently those terms occur in the full text.

\subsection{TF-IDF Cosine Similarity (Vector Space Model)}

The Vector Space Model represents both queries and documents as real-valued term weight vectors and ranks documents by the cosine of the angle between them \cite{salton1983}. The term weight scheme used is log-normalised TF-IDF:

\begin{equation}
\mathrm{idf}(t) = \log\!\left(\frac{N+1}{\mathrm{df}(t)+1}\right)+1
\end{equation}

\begin{equation}
w_d(t) = \bigl(1 + \log\,\mathrm{tf}(t,d)\bigr) \cdot \mathrm{idf}(t), \quad
w_q(t) = \bigl(1 + \log\,\mathrm{tf}(t,q)\bigr) \cdot \mathrm{idf}(t)
\end{equation}

\begin{equation}
\mathrm{sim}(q,d) = \frac{\displaystyle\sum_{t \in q \cap d} w_q(t)\,w_d(t)}{\lVert\mathbf{q}\rVert\,\lVert\mathbf{d}\rVert}
\end{equation}

The $+1$ smoothing in the IDF numerator and denominator prevents division-by-zero and avoids negative IDF for universal terms. Document L2-norms $\lVert\mathbf{d}\rVert$ are precomputed during the statistics phase and stored in \texttt{doc\_norms.json}, so query-time scoring requires only a dot-product accumulation over postings for query terms followed by a single division per candidate. This sparse scoring strategy is efficient because only documents containing at least one query term are ever scored.

\subsection{BM25}

BM25 is a probabilistic ranking function that extends TF-IDF with term-frequency saturation and explicit document-length normalisation \cite{robertson2009}:

\begin{equation}
\mathrm{idf}_{\mathrm{BM25}}(t) = \log\!\left(\frac{N - \mathrm{df}(t) + 0.5}{\mathrm{df}(t) + 0.5} + 1\right)
\end{equation}

\begin{equation}
\mathrm{score}(d,q) = \sum_{t \in q} \mathrm{idf}_{\mathrm{BM25}}(t) \cdot
\frac{\mathrm{tf}(t,d)\,(k_1+1)}{\mathrm{tf}(t,d) + k_1\!\left(1 - b + b\,\dfrac{|d|}{\mathrm{avgdl}}\right)}
\end{equation}

The parameters $k_1 = 1.2$ and $b = 0.75$ are the standard defaults recommended by Robertson and Zaragoza \cite{robertson2009}. The saturation mechanism ensures that term-frequency contribution asymptotically approaches a ceiling, preventing a single highly repeated term from dominating the score. The $b$ parameter controls the degree of length normalisation: at $b = 1$ documents are fully normalised by average length; at $b = 0$ no length normalisation is applied. The default $b = 0.75$ provides partial normalisation appropriate for heterogeneous corpora where document length varies substantially --- as it does in Project Gutenberg, where books range from short poems to multi-volume novels.

% ============================================================
\section{Evaluation}
% ============================================================

\subsection{Experimental Protocol}

The six assignment queries are hardcoded in the evaluator to ensure full reproducibility and eliminate any risk of accidental query drift between experiment runs:

\begin{enumerate}
    \item \emph{to be, or not to be}
    \item \emph{English Grammar}
    \item \emph{Philip K Dick}
    \item \emph{Jabberwocky}
    \item \emph{Gutenberg}
    \item \emph{Dornr\"{o}schen}
\end{enumerate}

For each query and each model, up to 100 results are exported to \texttt{<query\_nr>\_<model>.tsv} with columns: \texttt{rank}, \texttt{book\_id}, \texttt{score}, \texttt{preview}, \texttt{start\_line}. A total of 18 TSV files are produced ($6 \times 3$).

Because the assignment does not supply relevance judgements, standard effectiveness metrics such as Precision@K, MAP, and nDCG cannot be computed without making unverifiable assumptions \cite{manning2008}. Instead, evaluation uses three comparative metrics that characterise ranking behaviour across models on the same query set: \textbf{Intersection@K} (candidate-set agreement), \textbf{Ranking displacement@100} (ordering stability among shared candidates), and \textbf{Mean top-10 document length} (length-bias indicator).

\subsection{Aggregate Results}

Table~\ref{tab:agg} reports pairwise comparison metrics averaged across all six queries.

\begin{table}[t]
\caption{Aggregate pairwise comparison metrics (averaged over 6 queries).}
\label{tab:agg}
\centering
\begin{tabular}{lccc}
\toprule
Model Pair & Int@10 & Int@100 & Disp@100 \\
\midrule
Structured vs TF-IDF & 0.167 & 2.500  & 61.825 \\
Structured vs BM25   & 0.500 & 5.667  & 79.215 \\
TF-IDF vs BM25       & 2.833 & 31.500 & 34.404 \\
\bottomrule
\end{tabular}
\end{table}

Mean top-10 document lengths (tokens) were: Structured: 120,022; TF-IDF: 18,205; BM25: 95,414. The large TF-IDF/BM25 gap confirms that BM25's length normalisation promotes shorter documents more than log-normalised TF-IDF cosine similarity does. Structured retrieval returns long documents because metadata field length has no effect on its scoring function.

\subsection{Per-Query Results}

Table~\ref{tab:perquery} records per-query pairwise metrics produced directly by the evaluator, making the report auditable against \texttt{evaluation\_summary.txt}.

\begin{table*}[t]
\caption{Per-query pairwise comparison metrics (Int = Intersection, Disp = mean rank displacement).}
\label{tab:perquery}
\centering
\begin{tabular}{lrrrrrrrrr}
\toprule
Query & \multicolumn{3}{c}{Structured vs TF-IDF} & \multicolumn{3}{c}{Structured vs BM25} & \multicolumn{3}{c}{TF-IDF vs BM25} \\
\cmidrule(lr){2-4}\cmidrule(lr){5-7}\cmidrule(lr){8-10}
 & Int@10 & Int@100 & Disp & Int@10 & Int@100 & Disp & Int@10 & Int@100 & Disp \\
\midrule
to be, or not to be & 0 & 0  & 100.000 & 0 & 1  & 9.000   & 0 & 0  & 100.000 \\
English Grammar     & 1 & 6  & 33.000  & 3 & 10 & 22.900  & 0 & 30 & 25.367  \\
Philip K Dick       & 0 & 2  & 40.000  & 0 & 0  & 100.000 & 1 & 29 & 27.103  \\
Jabberwocky         & 0 & 0  & 100.000 & 0 & 0  & 100.000 & 8 & 43 & 4.140   \\
Gutenberg           & 0 & 5  & 10.200  & 0 & 24 & 55.292  & 0 & 10 & 35.400  \\
Dornr\"{o}schen     & 0 & 0  & 100.000 & 0 & 0  & 100.000 & 8 & 77 & 14.416  \\
\bottomrule
\end{tabular}
\end{table*}

\subsection{Discussion}

Four consistent patterns emerge from the results.

\textbf{Metadata-intent isolation.} Structured retrieval has near-zero overlap with both lexical models on most queries. Queries such as \emph{Jabberwocky} and \emph{Dornr\"{o}schen} are absent from most metadata fields, so structured retrieval returns empty or near-empty result lists. This is expected and correct: it reflects the fundamental difference in retrieval intent between a metadata channel and a full-text channel, not an implementation error.

\textbf{Lexical model coupling.} TF-IDF and BM25 share far more candidates than either shares with structured retrieval (Int@100 of 31.5 versus at most 5.7), which follows naturally from both models operating on the same term-frequency statistics from the same postings. Their non-trivial displacement (34.4) confirms that ranking order can still shift substantially as a result of BM25's saturation and length normalisation terms.

\textbf{Query-specific stability.} Highly distinctive, low-frequency terms (\emph{Jabberwocky}, \emph{Dornr\"{o}schen}) yield the strongest TF-IDF/BM25 top-10 agreement (Int@10 = 8). Common-term queries (\emph{to be, or not to be}) yield the weakest agreement, with Disp@100 = 100.0. This is consistent with the theoretical prediction that IDF weighting is most effective for rare terms and least effective for ubiquitous stopwords \cite{manning2008}.

\textbf{Length bias.} The mean top-10 document length gap between TF-IDF (18,205) and BM25 (95,414) is the most practically significant behavioural difference between the two lexical models. TF-IDF cosine normalisation divides by the full document norm, which penalises long documents uniformly across all terms and may over-correct for length on short, focused texts.

\subsection{Advantages and Disadvantages of Each Model}

\textbf{Structured retrieval} offers the highest interpretability and fastest query execution. It is the correct model for bibliographic lookup tasks. Its weakness is complete inability to retrieve thematically relevant documents whose metadata does not contain query terms.

\textbf{TF-IDF cosine} provides a strong, well-understood baseline for full-text lexical retrieval \cite{salton1983}. Cosine normalisation removes the document-length confound, making it relatively stable across documents of varying length. Its limitation is sensitivity to stopword-heavy queries and lexical mismatch.

\textbf{BM25} improves on TF-IDF through term-frequency saturation and tunable length normalisation and is the preferred default lexical ranker for heterogeneous corpora \cite{robertson2009}. Its limitation is that it remains a bag-of-words model incapable of capturing phrase structure, term proximity, or semantic similarity.

% ============================================================
\section{Reproducibility}
% ============================================================

\subsection{Validation Procedure}

Reproducibility was enforced through a five-step deterministic checklist: (1) compile all Python modules to detect syntax and import errors; (2) execute one query per model and verify non-empty ranked output; (3) run \texttt{run-experiments} using the hardcoded six-query list to produce all 18 TSV files; (4) verify each TSV for correct column headers; (5) inspect \texttt{evaluation\_summary.txt} for the presence of all pairwise metric rows.

\subsection{Core Commands}

\begin{lstlisting}
# Build index (with resume support)
python3 main.py --project-root . build-index \
    --batch-size 250 --resume

# Full assignment evaluation (hardcoded 6 queries)
python3 main.py --project-root . run-experiments \
    --output results/evaluation_summary.txt \
    --tsv-dir results

# Single interactive query
python3 main.py --project-root . search \
    --model bm25 --query "Jabberwocky" --top-k 10
\end{lstlisting}

% ============================================================
\section{Conclusions}
% ============================================================

This project delivers a complete IR pipeline from batch indexing to comparative evaluation over a large public-domain corpus. The central engineering decision --- a disk-backed batch indexer using SQLite partial indexes --- was essential for memory safety and build recoverability and would remain unchanged in any revision of the design.

The comparative evaluation demonstrates clear behavioural differences between the three retrieval approaches. Structured retrieval is the correct choice for bibliographic lookup but is orthogonal to full-text retrieval. TF-IDF and BM25 agree strongly on candidate sets but diverge in ranking order, particularly on stopword-heavy queries. BM25's length normalisation produces a notably different document-length distribution in top results, which is the most practically significant difference between the two lexical models on this corpus.

The primary limitation of the current evaluation is the absence of relevance judgements. Without qrels, direct effectiveness metrics cannot be computed. The most valuable extension would be to manually label top-10 results per query as relevant or non-relevant, enabling Precision@10 and MAP computation. Further improvements include phrase proximity scoring for exact-expression queries, language-aware normalisation for multilingual robustness, and incremental index updates to avoid full rebuilds when the corpus grows.

\begin{thebibliography}{9}

\bibitem{manning2008}
Christopher D. Manning, Prabhakar Raghavan, and Hinrich Sch\"{u}tze.
\newblock \textit{Introduction to Information Retrieval}.
\newblock Cambridge University Press, 2008.

\bibitem{robertson2009}
Stephen Robertson and Hugo Zaragoza.
\newblock The Probabilistic Relevance Framework: BM25 and Beyond.
\newblock \textit{Foundations and Trends in Information Retrieval}, 3(4):333--389, 2009.

\bibitem{salton1983}
Gerard Salton and Michael J. McGill.
\newblock \textit{Introduction to Modern Information Retrieval}.
\newblock McGraw-Hill, 1983.

\bibitem{zobel2006}
Justin Zobel and Alistair Moffat.
\newblock Inverted Files for Text Search Engines.
\newblock \textit{ACM Computing Surveys}, 38(2), 2006.

\end{thebibliography}

\end{document}